% -----------------------------------------------------------------------------
% This file is automatically generated.
% Do not edit manually.
% Source: notebooks/support/hypothesis-testing/support.Rmd
% -----------------------------------------------------------------------------


\noindent On page 120 we calculated significance levels for the occurrence of finding one
or two errors. These probabilities are calculated in Python with the following
code:

\par\addvspace{\topsep}
\begingroup
\fvset{listparameters={%
  \setlength{\topsep}{0pt}%
  \setlength{\partopsep}{0pt}%
  \setlength{\parsep}{0pt}%
  \setlength{\itemsep}{0pt}%
}}
\begin{Verbatim}[breaklines=true,formatcom=\color{ada_blue}]
hypergeom.cdf(1, 1200, 60, 45)
hypergeom.cdf(2, 1200, 60, 102)
\end{Verbatim}
\begin{Verbatim}[breaklines=true,formatcom=\color{ada_light_blue}]
[1] 0.3294032

[1] 0.09989878
\end{Verbatim}
\endgroup
\par\addvspace{\topsep}
\noindent Refer to Equation 2.1. Remember that Python's hypergeometric distribution
function uses the population size $N$, the number of deviations in the
population $M$, the sample size $n$, and the number of errors found $k$.
